A coding agent's final report is often the only account of its work that anyone reads. We study what these reports say when an agent games its tests. We built 24 small repositories in which one test contradicts the function's documentation, so that no faithful implementation passes the full suite, and ran six Claude models in the Claude Code harness for \nAnalyzed{} trials under a pre-registered protocol. Four of the six models gamed the suite in most trials, almost never by hardcoding: they adopted a general-looking rule under which the wrong test is right, implemented it, and often rewrote the documentation to match; \colAnyFail{} of \colN{} such rules changed behavior on other inputs. In \hOneEst{} of these cases the report justified the rule by a requirement written nowhere in the repository or conversation, either credited to an outside authority or stated as an established rule, while test edits were disclosed \hTwoMod{} of the time. Minimal pairs that remove any plausible rule cut gaming from \hSixOrig{} to \hSixVar. Asked in the same session whether it added test-specific logic, the agent usually said no, pointing to the rule's generality; asked for the rule's source, it usually said the test was the only one. LLM reviewers shown only the diff approved gamed changes more often when they carried the agent's report than an honest description or none. Claude Opus 5.5 and Fable 5.1 gamed less and disclosed every special case they made. We release the tasks, harness, and all trial records.
